\documentclass[12pt]{amsart} 

\input{./latex-preamble/cm-preamble.tex}

%\graphicspath{ {./diagrams/} }

\DeclareMathOperator{\pr}{pr}
\newcommand{\SPEC}{\text{\bf{Spec}}}
\newcommand{\sS}{\mathscr{S}}

\newtheorem{lemma}{Lemma}

\begin{document}

{Hartshorne} \hfill {3-08-2025}

\bigskip\bigskip

\begin{center}

{\sc Chapter II.5, Week 7}

\end{center}

\begin{center}

  {Noah Parker} %\footnote{Collaborated with Kalev Martinson and Frank :).}}
    
\end{center}

\bigskip\bigskip

\begin{enumerate}
  \item[5.9.] {
      Let $S$ be a graded ring, generated by $S_1$ as an $S_0$-algebra, $M$ a graded $S$-module, and $X=\Proj S$.
      \begin{enumerate}
        \item Show there is a natural homomorphism $\alpha:M\to \Gamma_*(\wtilde{M})$.
        \item Assume now that $S_0=A$ is a finitely generated $k$-algebra for some field $k$, that $S_1$ is a finitely generated $A$-module, and that $M$ is a finitely generated $S$-module.
          Show that the map $\alpha$ is an isomorphism in all large enough degrees, i.e., there is a $d_0\in \ZZ$ such that for all $d\geqs d_0$, $\alpha_d:M_d\to \Gamma(X,\wtilde{M}(d))$ is an isomorphism.
          [\emph{Hint}: Use the methods of the proof of (5.19).]
        \item With the same hypothesis, we define an equivalence relation $\approx$ on graded $S$-modules by saying $M\approx M'$ if there is an integer $d$ such that $M_{\geqs d}\cong M_{\geqs d}'$.
          Here, $M_{\geqs d}=\bigoplus_{n\geqs d}M_n$.
          We say that a graded $S$-module is \emph{quasi-finitely generated} if it is equivalent to a finitely generated module.
          Now show that the functors $\wtilde\cdot$ and $\Gamma_*$ induces an equivalence of categoriese between the category of quasi-finitely generated graded $S$-modules modulo the equivalence relation $\approx$, and the category of coherent $\O_X$-modules.
      \end{enumerate}
    }
\end{enumerate}

\begin{proof}
  (a) Fix generators $\{x_i\}_{i\in I}\subset S_1$ of $S$ as an $S_0$-algebra.
  This gives us an affine open cover $\{D_+(x_i)\cong \Spec S_{(x_i)}\}_{i\in I}$ of $\Proj S$ such that 
  \begin{align*}
    \wtilde{M}|_{D_+(x_i)}(n)
    &\cong \wtilde{M(n)}|_{D_+(x_i)} \\
    &\cong (M(n)_{(x_i)})^\sim \\
    &\cong ((M_{x_i})_n)^\sim
  \end{align*}
  since restriction commute with tensor products, and $x_i\in S_1$ implies 
  \[
    (M\otimes_S N)_{(x_i)}=M_{(x_i)}\otimes_{S_(x_i)}N_{(x_i)}.
  \]
  Then
  \begin{align*}
    \Gamma_*(\wtilde{M}|_{D_+(x_i)})
    &\cong \bigoplus_{n\in \ZZ} \Gamma(D_+(x_i), \wtilde{M}|_{D_+(x_i)}(n)) \\
    %&\cong \bigoplus_{n\in \ZZ} \Gamma(D_+(x_i), \wtilde{M(n)}|_{D_+(x_i)}) \\
    &\cong \bigoplus_{n\in \ZZ} \Gamma(\Spec S_{(x_i)}, ((M_{x_i})_n)^\sim) \\
    &\cong \bigoplus_{n\in \ZZ}(M_{x_i})_n \\
    &\cong M_{x_i}
  \end{align*}
  supplies us with a natural map $\alpha_i:M\hookrightarrow M_{x_i}\iso \Gamma_*(\wtilde{M}|_{D_+(x_i)})$ for each $i$.
  Using a similar argument to the above, we define maps $\alpha_{ij}:M\hookrightarrow M_{x_ix_j}\iso \Gamma_+(\wtilde{M}|_{D_+(x_ix_j)})$ and show that $\rho_{D_+(x_ix_j)}\circ \alpha_i= \alpha_{ij} =\rho_{D_+(x_ix_j)} \circ \alpha_j$.

  Fixing $m\in M$, the $\alpha_i(m)_n\in \Gamma(D_+(x_i),\wtilde{M}|_{D_+(x_i}(n))$ glue together to a section $\alpha(m)_n\in \Gamma(X,\wtilde{M})$.
  Only finitely many $\alpha_i(m)_n$ are nonzero for each $i$, so only finitely many $\alpha(m)_n$ are nonzero, whence $\alpha(m)=\sum_n\alpha(m)_n\in \Gamma_*(\wtilde{M})$ is well-defined.
  We take $\alpha:m\mapsto \alpha(m)$ to be our natural map.

  (b) 
\end{proof}

\begin{enumerate}
  \item[5.10.] {
      Let $A$ be a ring, $S=A[x_0,\ldots, x_r]$, and $X=\Proj S$.
      We have seen that a homogeneous ideal $I$ in $S$ defines a closed subscheme of $X$ (Ex. 3.12), and that convrsely every closed subscheme of $X$ arises in this way (5.16).
      \begin{enumerate}
        \item For any homogeneous ideal $I\subset S$, we define the saturation $\ol{I}$ of $I$ to be
          \[
            \ol{I}:=\{s\in S\mid \forall 0\leqs i\leqs r~(\exists n>0~(x_i^ns\in I))\}
          \] 
          We say that $I$ is \emph{saturated} if $I=\ol{I}$.
          Show that $\ol{I}$ is a homogeneous ideal of $S$.
        \item Two homogeneous ideals $I_1$ and $I_2$ of $S$ define the same closed subscheme of $X$ iff they have the same saturation.
        \item If $Y$ is any closed subscheme of $X$, then the ideal $\Gamma_*(\I_Y)$ is saturated.
          Hence, it is the largest homogeneous ideal defining the subscheme $Y$.
        \item There is a 1-1 correspondence between saturated ideals of $S$ and closed sub-schemes of $X$.
      \end{enumerate}
    }
\end{enumerate}

\begin{proof}
  (a) If $s\in \ol{I}$, then we can write $s= \sum_ds_d$ as the sum of its homogeneous parts, so that $x_i^ns=\sum_dx_i^ns_d\in I$ and $I$ homogeneous implies $x_i^ns_d\in I$ for each $d$.
  We conclude that $s_d\in \ol{I}$ for all $d$, whence $\ol{I}$ is homogeneous.

  (b) Fix homogeneous ideals $I_1,I_2\lideal S$. 
  We have an affine cover $D_+(x_i)\cong \Spec S_{(x_i)}$ of $X$, so it suffices to show that 
  \[
    \frac{S_{(x_i)}}{(I_1)_{(x_i)}}
    \cong \wtilde{S/I_1}|_{D_+(x_i)}
    \cong \wtilde{S/I_2}|_{D_+(x_i)} 
    \cong \frac{S_{(x_i)}}{(I_2)_{(x_i)}}
  \] 
  for all $0\leqs i\leqs r$ iff $\ol{I}_1=\ol I_2$.
  This occurs iff the kernels of the respective projection maps are the same, i.e. $(I_1)_{(x_i)}= (I_2)_{(x_i)}$.
  %Clearly this holds if $(I_1)_{x_i}=(I_2)_{x_i}$, so suppose conversely that the equality holds only in degree 0.
  %It suffices to show that $
  If $\ol I_1=\ol I_2$, then $\exists n>0$ with $x_i^ns\in I_1$ iff $\exists n>0$ with $x_i^ns\in I_2$.

  (c) 
\end{proof}

\begin{enumerate}
  \item[5.11.] {
      Let $S$ and $T$ be two graded rings with $S_0=T_0=A$.
      We define the \emph{Cartesian product} $S\times_A T$ to be the graded ring $\bigoplus_{d\geqs 0}S_d\otimes_A T_d.$
      If $X=\Proj S$ and $Y=\Proj T$, show that $\Proj (S\times_A T) \cong X\times_A Y$, and show that the sheaf $\O(1)$ on $\Proj(S\times_A T)$ is isomorphic to the sheaf $p_1^*(\O_X(1))\otimes p_2^*(\O_Y(1))$ on $X\times_A Y$.

      The Cartesian product of rings is related to the Segre embedding of projective spaces in the following way.
      If $x_0,\ldots, x_r$ is a set of generated for $S_1$ over $A$, corresponding to a projective embedding $X\hookrightarrow \PP_A^r$, and if $y_0,\ldots, y_s$ is a set of generators for $T_1$, corresponding to a projective embedding $Y\hookrightarrow \PP_A^s$, then $x_i\otimes y_j$ is a set of generators for $(S\times_A T)_1$, and hence defines a projective embedding $\Proj (S\times_A T)\hookrightarrow \PP_A^N$, with $N=rs + r + s-1$
      This is just the image of $X\times Y\subset \PP^r\times \PP^s$ in its Sege embedding.
    }
\end{enumerate}

\begin{proof}
\end{proof}

\begin{enumerate}
  \item[5.12.] {
      \begin{enumerate}
        \item Let $X$ be a scheme over a scheme $Y$, and let $\L$, $\M$ be two very ample invertible sheaves on $X$.
          Show that $\L\otimes \M$ is also very ample.
          [\emph{Hint}: Use a Segre embedding.]
        \item Let $f:X\to Y$ and $g:Y\to Z$ be two morphisms of schemes.
          Let $\L$ be a very ample invertible sheaf on $X$ relative to $Y$, and let $\M$ be a very ample invertible sheaf on $Y$ relative to $Z$.
          Show that $\L\otimes f^*\M$ is a very ample invertible sheaf on $X$ relative to $Z$.
      \end{enumerate}
    }
\end{enumerate}

\begin{proof}
\end{proof}

\begin{enumerate}
  \item[5.13.] {
      Let $S$ be a graded ring, generated by $S_1$ as an $S_0$-algebra.
      For any integer $d>0$, let $S^{(d)}$ be the graded ring $\bigoplus_{n\geqs 0}S_n^{(d)}$, where $S_n^(d)=S_{nd}$.
      Let $X=\Proj S$.
      Show that $\Proj S^{(d)}\cong X$, and that the sheaf $\O(1)$ on $\Proj S^{(d)}$ corresponds via this isomorphism to $\O_X(d)$.
      
      This construction is related to the $d$-uple \emph{embedding} in the following way.
      Denote $A=S_0$.
      If $x_0,\ldots, x_r$ is a set of generators for $S_1$, corresponding to an embedding $X\hookrightarrow \PP_A^r$, then the set of monomials of degree $d$ in the $x_i$ is a set of generators for $S_1^{(d)}=D_d$.
      These define a projective embedding of $\Proj S^{(d)}$ which is none other than the image of $X$ under the $d$-uple embedding of $\PP_A^r$.
    }
\end{enumerate}

\begin{enumerate}
  \item[5.14.] {
      Let $A$ be a ring, and let $X$ be a closed subscheme of $\PP_A^r$.
      We define the homogeneous coordinate ring $S(X)$ of $X$ for the given embedding to be $A[x_0,\ldots, x_r]/I$, where $I$ is the ideal $\Gamma_*(\I_X)$ constructed in the proof of (5.14).
      (Of course, if $A$ is a field and $X$ a variety, this coincides with the usual definition!)
      Recall that a scheme $X$ is \emph{normal} if its local rings are integrall closed domains.
      A closed subscheme $X\subset \PP_A^r$ is \emph{projectively normal} for the given embedding, if the homogeneous coordinate ring $S(X)$ is an integrally closed domain .
      Now assume that $A$ is a finitely generated $k$-algebra for some field $k$, and that $X$ is a connected, normal closed susbcheme of $\PP_A^r$.
      Show that for some $d>0$, the $d$-uple embedding of $X$ is projectively normal, as follows.
      \begin{enumerate}
        \item Let $S$ be the homogeneous coordinate ring of $X$, and let $S'=\bigoplus_{n\geqs 0}\Gamma(X,\O_X(n))$.
          Show that $S$ is a domain, and that $S'$ is its integral closure.
          [\emph{Hint}: First show that $X$ is integral.
          Then regard $S'$ as the global sections of the sheaf of rings $\sS=\bigoplus_{n\geqs 0}\O_X(n)$ on $X$, and show that $\sS$ is a sheaf of integrally closed domains.]
        \item Use (Ex. 5.9) to show that $S_d=S_d'$ for all sufficiently large $d$.
        \item Show that $S^{(d)}$ is integrally closed for sufficiently large $d$, and hence conclude that the $d$-uple embedding of $X$ is projectively normal.
        \item As a corolary of (a), show that a closed subscheme $X\subset \PP_A^r$ is projectively normal if and only if it is normal, and for every $n\geqs 0$ the natural map $\Gamma(\PP^r,\O_{\PP^r}(n))\to \Gamma(X,\O_X(n))$ is surjective.a
      \end{enumerate}
    }
\end{enumerate}

\end{document}
